\section{Metodologia}\label{sec:metodo}

\subsection{Desenho e políticas comparadas}

Adota-se uma avaliação quantitativa e comparativa por simulação de eventos discretos \cite{banks2010}: as políticas são executadas com os mesmos pedidos, frota, jornadas, compromissos, método de distância e parâmetros em cada comparação, e suas diferenças são calculadas por período. Nos dados históricos, o desenho é contrafactual e retrospectivo, pois reaplica decisões alternativas a entradas registradas; nos cenários sintéticos, varia condições geradas para examinar o método. A diferença estimada é interna ao simulador; o despacho imediato é uma regra de referência, não uma reprodução do despacho praticado pela empresa. O Quadro~\ref{qua:politicas} define as quatro políticas.

Distinguem-se três verificações. Testes automáticos examinam se a implementação segue a formulação. A validação de componentes confronta distâncias, velocidades, durações e disponibilidade com registros observados; foi parcial e posterior aos resultados principais (\nameref{sec:validade}). A validação operacional exigiria reproduzir o comportamento da operação; esse objetivo não foi alcançado.

\begin{quadro}[!htb]
\caption{Políticas comparadas: despacho imediato (A), lote sequencial (B), atribuição conjunta (C) e guloso global (D)}\label{qua:politicas}
\centering\fontetabela
\begin{tabular}{|>{\raggedright\arraybackslash}p{34mm}|>{\raggedright\arraybackslash}p{60mm}|>{\raggedright\arraybackslash}p{54mm}|}
\hline
\textbf{Política} & \textbf{Quando decide} & \textbf{Como escolhe} \\
\hline
despacho imediato & na chegada do pedido (fila por ordem de chegada se não houver veículo elegível) & veículo elegível mais próximo \\
\hline
lote sequencial & a cada 60 s; pedidos em ordem de chegada e, em empate, por identificador & cada pedido recebe o veículo elegível mais próximo ainda livre \\
\hline
atribuição conjunta & a cada 60 s & solução da equação~\ref{eq:obj} pelo Húngaro \\
\hline
guloso global$^\dagger$ & a cada 60 s & pares elegíveis do mais curto para o mais longo, sem repetir pedido nem veículo \\
\hline
\end{tabular}
\fonte{Elaboração própria. A e B: referências; C: conjunta; D: análise posterior.}
\end{quadro}

A comparação principal, fixada antes das execuções, é atribuição conjunta contra despacho imediato: a troca completa de política. A comparação com o lote sequencial isola a decisão conjunta, porque as duas esperam o mesmo lote. A comparação entre lote sequencial e despacho imediato mostra o efeito de esperar o lote mantendo a escolha sequencial. A ordem do lote sequencial é a mesma do despacho imediato e foi definida antes de qualquer resultado.

\subsection{Despacho em lotes}

Os lotes ocorrem a cada 60 s, contados da meia-noite. Um pedido espera o próximo lote, e sua espera é contada desde a chegada. Um veículo liberado entre dois lotes só volta a ser candidato no lote seguinte. No despacho imediato, o veículo liberado pode receber imediatamente o pedido pendente mais antigo para o qual é elegível. A cada lote, o simulador constrói a lista de pares elegíveis, aplica a política e rejeita pares inelegíveis ou repetidos. A Figura~\ref{fig:fluxo} resume o procedimento; o Apêndice~\ref{ap:metodo} detalha a matriz de custos.

\begin{figure}[!htb]
\caption{Fluxo da atribuição em lote}\label{fig:fluxo}
\centering
\begin{tikzpicture}[node distance=7mm, box/.style={draw, rounded corners=2pt, align=center, font=\fontsize{9}{10}\selectfont, minimum height=10mm, text width=34mm}, >=Stealth]
\node[box](a){Instante de lote $t$ (múltiplo de 60 s)};
\node[box,right=of a](b){Pedidos pendentes (chegada $\le t$)};
\node[box,right=of b](c){Veículos elegíveis (livres, no raio, sem conflito)};
\node[box,below=of c](d){Matriz de custos $c_{rv}$ (km)};
\node[box,left=of d](e){Escolha: lote sequencial ou atribuição conjunta};
\node[box,left=of e](f){Validação: pares únicos e elegíveis};
\node[box,below=of f](g){Atualização de ocupação, posição e liberação};
\node[box,right=of g](h){Pedidos sem par seguem para $t+60$ s};
\draw[->](a)--(b);\draw[->](b)--(c);\draw[->](c)--(d);\draw[->](d)--(e);\draw[->](e)--(f);\draw[->](f)--(g);\draw[->](g)--(h);
\end{tikzpicture}
\par\fonte{elaborada pelos autores a partir do simulador.}
\end{figure}

\subsection{Dados reais}

Os dados vêm de uma empresa de transporte por aplicativo que opera na Grande Vitória (ES), restritos a um produto, à região da Grande Vitória e a quatro tipos de serviço, sem os serviços de categoria superior (material suplementar). Foram extraídos em modo somente leitura do banco operacional e do índice de posições dos veículos: pedidos (origem, destino, horário, tipo), jornadas de motoristas, compromissos externos, vínculos entre veículo e motorista, posições dos veículos (a cada cerca de 2 minutos) e cadastro de veículos. Nenhuma quilometragem do estudo vem de odômetro. A Tabela~\ref{tab:periodos} traz os períodos. Treino e validação serviram ao desenvolvimento. Setembro (replicação) e outubro (avaliação final) nunca tinham sido simulados antes da execução congelada.

\begin{table}[!htb]
\caption{Períodos e volumes dos dados reais (pedidos de tempo real, sem agendados)}\label{tab:periodos}
\centering\fontetabela
\begin{tabular}{|l|l|r|r|r|r|}
\hline
\textbf{Período} & \textbf{Datas (2026)} & \textbf{Dias} & \textbf{Pedidos} & \textbf{Motoristas/dia (média)} & \textbf{Compromissos} \\
\hline
Treino & 01/02--15/03 & 43 & 51.909 & 106,8 & 1.458 \\
\hline
Validação & 16/03--31/03 & 16 & 20.178 & 111,3 & 675 \\
\hline
Replicação & 16/09--30/09 & 15 & 16.918 & 86,7 & 674 \\
\hline
Final & 01/10--25/10 & 25 & 23.523 & 74,3 & 842 \\
\hline
\end{tabular}
\fonte{elaborada pelos autores a partir dos dados da operação. Pedidos agendados ficam fora; pedidos cancelados permanecem como demanda.}
\end{table}

As regras de construção foram fixadas antes da execução congelada e detalhadas no Apêndice~\ref{ap:metodo}. Um dia exige ao menos 300 instantes de posição e 200 pedidos; todos os 15 dias de setembro e os 25 de outubro cumpriram a regra. A disponibilidade é reconstruída retrospectivamente pela interseção da jornada realizada com a presença no registro de posições. Depois da primeira posição observada, a ocupação e a posição de cada veículo evoluem conforme as corridas atribuídas pelo simulador.

A duração com passageiro é a observada ou, se ausente, a mediana do dia calculada retrospectivamente. A política não recebe essa duração como campo, mas a elegibilidade fornecida pelo avaliador pode depender dela ao verificar conflitos com compromissos. Trata-se de informação retrospectiva na construção do ambiente, igual para as políticas, e não de informação que estaria necessariamente disponível em uma operação em tempo real. Novas tentativas do mesmo passageiro não foram deduplicadas, e cancelamentos posteriores à atribuição não são modelados. A queda conjunta no número de motoristas e pedidos entre treino e outubro pode refletir mudança operacional ou de registro; o estudo não distingue essas causas.

A aproximação histórica é estimada pela distância em linha reta multiplicada pelo fator rua/reta de 1,3012, com velocidade de 16,079 km/h. O raio é de 4,5 km em linha reta e a paciência, de 180 s. Os parâmetros vêm da calibração de março e permaneceram fixos em setembro e outubro; a recalibração posterior é descrita na seção \nameref{sec:validade}.

Identificadores de motoristas, veículos e passageiros foram pseudonimizados por HMAC-SHA256 com segredo fora do repositório; o artigo apresenta apenas agregados. Código, configurações, manifestos de verificação e resultados agregados documentam a reprodução; os cenários sintéticos usam geografia pública, mas a avaliação histórica depende de acesso autorizado às entradas confidenciais. Essas medidas não constituem, por si, uma avaliação de conformidade com a legislação de proteção de dados \cite{brasil2018}.

\subsection{Cenários sintéticos}

Os cenários sintéticos usam a geografia real de Vitória, Vila Velha, Serra e Cariacica: malhas do IBGE \cite{ibgemalhas} e rede viária do OpenStreetMap \cite{osm}, obtidas em 08/11/2026. As distâncias são caminhos mínimos pela rede (Dijkstra) entre 368 pontos de parada, escolhidos com células H3 \cite{uberh3}. A demanda, a frota, as jornadas e a velocidade são hipóteses do gerador, não calibradas. A frota tem 100 veículos, com jornadas de 9 horas, velocidade de 25 km/h, raio de 8 km por via e paciência de 600 s. Como esses parâmetros diferem dos dados reais, os percentuais das duas avaliações não são comparáveis em nível.

Há cinco cenários principais: demanda uniforme, demanda concentrada, frota inicialmente desequilibrada, mudança inesperada da demanda e demanda que migra de região. Outros 17 são variantes; nove alteram apenas o sorteio e não representam novas condições operacionais. Cada cenário foi simulado com 40 sementes reservadas, que são as unidades de análise. Como os resultados sintéticos foram vistos antes da decisão de adotar o método e os cenários compartilham rede, frota e gerador, sua evidência é exploratória e complementar. O Apêndice~\ref{ap:metodo} detalha a geografia e as variantes.

\subsection{Congelamento e origem da janela de 60 s}

A janela de 60 s foi escolhida no treino (01/02--15/03) por uma regra escrita antes da primeira execução: entre as alternativas com atendimento preservado e intervalo de km abaixo de zero, selecionar a de menor km, com desempate por simplicidade e tempo. A grade incluiu janelas de 30, 60 e 120 s. A janela de 120 s teve maior redução de km frente ao despacho imediato, mas atendeu 416 pedidos a menos e não cumpriu o critério de atendimento. A validação (16--31/03) e os cenários sintéticos foram vistos antes da decisão de adotar o método, embora não tenham escolhido a janela nem o algoritmo.

Em 09/11/2026, antes de simular setembro e outubro, foram congelados localmente código, parâmetros, políticas, comparações e critérios; o programa recusa a execução se o código mudar. Setembro foi executado antes de outubro, sem ajuste entre eles, falhas técnicas ou reexecuções. Esses períodos ficaram fora da seleção de janela e método, mas compartilham empresa, região, simulador e parâmetros de março. Os arquivos e hashes locais não constituem pré-registro externo verificável. A janela posterior de 26/10 a 08/11, reservada para a recalibração com janeiro a agosto, é o único teste denominado fora da amostra.

\subsection{Métricas e estatística}

O km vazio de uma política em um período soma as distâncias de aproximação dos pedidos despachados; não inclui reposicionamento nem retorno. A diferença relativa é
\begin{equation}
\Delta\% = 100\,\frac{K_C-K_{ref}}{K_{ref}}, \label{eq:delta}
\end{equation}
em que $K_C$ é o km vazio da atribuição conjunta e $K_{ref}$, o da referência; valor negativo significa redução. \emph{Pedido elegível} é aquele incluído no recorte de demanda do dia, antes da decisão de despacho. Nas saídas do simulador, \emph{atendido} é o pedido atribuído a um veículo no evento de despacho; a contagem não equivale a corridas concluídas na operação. $\Delta N$ é a diferença entre essas contagens, e a diferença de taxas usa todos os pedidos elegíveis como denominador. O \emph{aceite observado}, empregado na validação, exige registro de motorista a caminho e pertence a outra população e outro processo.

A espera simulada vai da chegada do pedido ao encontro estimado com o veículo, inclusive o tempo até o fechamento do lote. Ela é calculada apenas entre os atendidos de cada política. Como esses conjuntos podem diferir, a comparação de espera é descritiva.

O critério de atendimento, definido na introdução, foi fixado antes, sem acordo de nível de serviço com a empresa. Uma redução só é classificada como \enquote{com atendimento preservado} se o intervalo de 95\% da diferença ficar abaixo de zero e o critério de atendimento for cumprido.

A unidade de análise é o dia. Para cada dia, calcula-se a diferença de km entre as duas políticas. O intervalo de confiança de 95\% (IC95\%) indica a faixa de valores da diferença média diária compatível com a variação entre dias; ele é obtido por \textit{bootstrap}, isto é, sorteando blocos de 7 dias consecutivos muitas vezes e recalculando a média. Com 15 e 25 dias há poucos blocos, e o intervalo é descritivo; o material suplementar mostra intervalos alternativos. Um intervalo que inclui zero não mostra que as políticas são equivalentes: isso exigiria uma margem fixada antes, e nenhuma foi fixada. Nos cenários sintéticos, a unidade é a semente.

Após as referências deste PDF, o Apêndice~\ref{ap:metodo} detalha método e dados; o Apêndice~\ref{ap:repro} reúne verificações e reprodutibilidade. O PDF \caminho{article/supplement/material_suplementar.pdf} contém os apêndices suplementares A (método e dados), B (verificações e análises posteriores), C (resultados sintéticos por cenário) e D (rastreabilidade técnica), com paginação própria.
